\section{Proofs}
\label{app:proofs}

\paragraph{Proposition~\ref{prop:free}.}
$\mathbb{N}^P$ is the free commutative monoid on $\{1,\dots,P\}$: every $m$ is uniquely $\sum_j m_je_j$. A homomorphism
$F$ with $F(e_j)=s(j)$ must therefore satisfy $F(m)=\sum_jm_jF(e_j)=\sum_jm_js(j)$, and this formula is a
homomorphism, because addition in $G$ is commutative and associative~\citep{maclane1998categories}.

For the converse, let $F_0=F(\emptyset)$ be the prior's natural parameters and set $s(j)=F(\{j\})-F_0$. If
$F(A\sqcup B)=F(A)+F(B)-F_0$ for all disjoint $A,B$, induction on $|S|$ gives $F(S)=F_0+\sum_{j\in S}s(j)$. The
posterior is thus the prior times one factor per observed sensor. \hfill$\square$

\paragraph{Proposition~\ref{prop:pool}.}
Both mixture components are non-negative. Hence
$p_{\mathrm{pool}}(y)\geq\pi\,p_{\mathrm{site}}(y)$ and $p_{\mathrm{pool}}(y)\geq(1-\pi)\,p_{\mathrm{free}}(y)$.
Taking $-\log$ of each gives the two bounds.

Marginalization commutes with the linear pool by linearity of the integral:
$\int(\pi p_1+(1-\pi)p_2)\,dh=\pi\int p_1\,dh+(1-\pi)\int p_2\,dh$. McConway~\citep{mcconway1981marginalization} shows
that, under mild conditions, linear pools are the only pooling operators with this property. \hfill$\square$

\paragraph{Proposition~\ref{prop:init}.}
With a zero site head, every corrected site equals its anchor site. The lifted posterior is then exactly the FA-Gaussian
posterior of the anchor (Proposition~\ref{prop:exact} applied to the fitted factor model). With gate bias $+4$ and zero
gate weights, $\pi=\sigma(4)$ for every query. Apply Proposition~\ref{prop:pool}:
$-\log\sigma(4)=\log(1+e^{-4})=0.0180$. \hfill$\square$

\paragraph{Scale equivariance of SF.}
Let $y'=ay+b$ with $a>0$ be applied to the support and query targets of an episode. The SF pipeline standardizes by
the support mean and standard deviation. These transform as $m'=am+b$ and $s'=as$ (the floor
$10^{-4}(1+|m|)$ is inactive whenever the support target varies). The standardized targets are therefore unchanged.
The inputs do not depend on the target's units, so the network and the refitted anchor produce the same standardized
predictive $\tilde q$. The raw-unit predictive is $q'(y')=\tilde q((y'-m')/s')/s'$, so
$\NLL'=\NLL+\log a$, $\mathrm{CRPS}'=a\,\mathrm{CRPS}$ and the squared error scales by $a^2$. The test suite checks this
for $a=10^{-3}$ to within $2\times10^{-3}$ nats, including the 8-member test-time ensemble. \hfill$\square$

\paragraph{Proposition~\ref{prop:scalar}.}
Under the model, $p(x_j\mid y)=\N(x_j;a_jy,c_j)$ with $c_j=d_j^2+\sigma_j^2$.

If $q(y\mid x_j)=p(y\mid x_j)$ for all $x_j$, then $p(y)\,t_j(y;x_j)\propto p(y)\,p(x_j\mid y)$ as functions of $y$.
Hence $t_j(y;x_j)=\kappa_j(x_j)\,p(x_j\mid y)$, which is Gaussian in $y$ with precision $a_j^2/c_j>0$. For a pair, $q$
is therefore the naive-Bayes posterior $p(y)\,p(x_j\mid y)\,p(x_k\mid y)$. Its posterior mean is linear in
$(x_j,x_k)$ with coefficients
\[
  (a_j/c_j,\;a_k/c_k)\,/\,\pi_q, \qquad \pi_q=1+a_j^2/c_j+a_k^2/c_k ,
\]
so the coefficient of $x_j$ has the sign of $a_j$.

The Bayes posterior uses $C=\mathrm{Cov}(x_j,x_k\mid y)$, with off-diagonal entry $d_jd_k$. Its mean coefficients are
$(C^{-1}a)^\top/\pi$, with $\pi=1+a^\top C^{-1}a$. Both posteriors are Gaussian, so they coincide for all $(x_j,x_k)$
if and only if $C^{-1}a=\mathrm{diag}(C)^{-1}a$. Equivalently $a=C\,\mathrm{diag}(C)^{-1}a$, whose off-diagonal
components are $d_jd_k\,a_k/c_k=0$ and $d_jd_k\,a_j/c_j=0$. With $a_ja_k\neq0$ this holds if and only if $d_jd_k=0$.

For the sign flip, take $a=(1,1)$, $d=(1,0.5)$ and $\sigma=(0.05,0.05)$. Then the Bayes coefficient of $x_1$ is
$+0.499$ alone and $-0.959$ in the pair. \hfill$\square$

\paragraph{Proposition~\ref{prop:exact}, exactness.}
Given $z=(y,h)$, the readings are independent: $x_j\mid z\sim\N(w_j^\top z+c_j,\psi_j)$. Hence
\[
  p(z\mid x_S)\propto p(z)\prod_{j\in S}\N(x_j;w_j^\top z+c_j,\psi_j).
\]
In natural parameters each factor contributes $w_jw_j^\top/\psi_j$ to the precision and $w_j(x_j-c_j)/\psi_j$ to the
shift, which is \eqref{eq:site} with $o_j=x_j-c_j$. The prior is Gaussian, so the product is the Gaussian
\eqref{eq:site}. This holds for every $S$ because the factors of unobserved sensors are absent from both sides.
Marginalizing $h$ gives $p(y\mid x_S)$. \hfill$\square$

\paragraph{Proposition~\ref{prop:rank}, minimal lifting.}
\emph{Step 1: the product is a latent-factor model.} By the computation in the proof of Proposition~\ref{prop:exact}
read backwards, a product of sites \eqref{eq:site} with $o_j=x_j-c_j$ is the exact posterior over $z$ of the model
$M_q$: $z\sim\N(\mu_z,\Sigma_z)$ and, independently given $z$, $x_j\sim\N(w_j^\top z+c_j,\psi_j)$. Its $y$-marginal is
$M_q$'s conditional $p_q(y\mid x_S)$, and $M_q$'s sensor covariance is $\Sigma_q=W\Sigma_zW^\top+\mathrm{diag}(\psi)$
with $W\in\R^{P\times(1+K')}$, so $\Sigma_q-\mathrm{diag}(\psi)$ has rank at most $1+K'$.

\emph{Step 2: conditionals on $|S|\leq2$ identify the covariance.} Write $V=\mathrm{Var}(y)$, $c=\mathrm{Cov}(x,y)$
and $\Sigma=\mathrm{Cov}(x)$. The case $S=\emptyset$ gives $V$. For $S=\{j\}$ the conditional has slope
$\beta_j=c_j/\Sigma_{jj}$ and variance $r_j=V-c_j\beta_j$, and $\beta_j\neq0$ because $c_j=a_j\neq0$. Hence
$c_j=(V-r_j)/\beta_j$ and $\Sigma_{jj}=c_j/\beta_j$. For $S=\{j,k\}$ the slope vector $b\neq0$ solves
$\Sigma_{\{jk\}}b=(c_j,c_k)$, and whichever entry of $b$ is non-zero determines $\Sigma_{jk}$ from one row. Since
$M_q$ has the same conditionals, the same formulas apply to it, so $\Sigma_q=\Sigma=LL^\top+\mathrm{diag}(\sigma^2)$.

\emph{Step 3: fewer factors cannot reproduce $\Sigma$.} Suppose toward contradiction $K'<K$. Let $r=1+K$, so $LL^\top$ has rank $r$. Suppose
$\Sigma=G+\Psi'$ with $\Psi'$ diagonal and $\mathrm{rank}(G)\leq r$. Fix a row $i$ and two disjoint row sets $I_1,I_2$
not containing $i$ with $L_{I_1},L_{I_2}$ nonsingular. In the $(r+1)\times(r+1)$ submatrix of $\Sigma$ with rows
$\{i\}\cup I_1$ and columns $\{i\}\cup I_2$, every entry except $(i,i)$ lies off the diagonal, so it is shared by $G$
and $LL^\top$. Both $G$ and $LL^\top$ have rank at most $r$, so the corresponding minors vanish. Expanding along entry
$(i,i)$, its cofactor is $\det(L_{I_1}L_{I_2}^\top)\neq0$, so $G_{ii}=(LL^\top)_{ii}$. Hence $G=LL^\top$, which has
rank $r$~\citep{anderson1956statistical}. Applying this to $G=W\Sigma_zW^\top$ from Step 1 gives $1+K'\geq r=1+K$.

For almost every $L$ all $r\times r$ submatrices are nonsingular, and $P\geq2r+1=2K+3$ leaves two disjoint sets of $r$
rows after deleting any row. \hfill$\square$



\section{Protocols, endpoints and amendments}
\label{app:protocols}

\paragraph{Pre-registration.} Every protocol was committed to the public repository before any score on its test data
existed. The repository history is available to reviewers. Each protocol fixes:
\begin{itemize}
  \item the networks, splits and test panels (seeded);
  \item the models and recipes;
  \item the endpoints, and the statistic and pass rule of each.
\end{itemize}

Amendments were committed before the scores they affect. Table~\ref{tab:all-endpoints} lists every endpoint and its
outcome. The amendments were:
\begin{itemize}
  \item \textbf{v6, amendments 2 and 3:} undefined closed-form and TabPFN scores on degenerate episodes are recorded as NaN.
  \item \textbf{v8:} a vectorized but distributionally identical kriging sampler.
  \item \textbf{v9, amendment 1:} PFN-L was retrained, because the v4 checkpoints had not been kept.
  \item \textbf{v10, amendment 1:} a time-zone bug in the USGS parser was fixed before any panel was built.
\end{itemize}

\paragraph{Statistics.} Endpoints compare per-episode metrics:
\begin{itemize}
  \item learned models are averaged over three seeds;
  \item episodes are pooled over networks;
  \item paired 95\% intervals are cluster-robust over network-time-windows;
  \item NLL endpoints pass when the gain is at least $0.01$ nats and the lower bound is above zero.
\end{itemize}

\begin{table}[h]\centering\scriptsize
\caption{\textbf{Every pre-registered endpoint} of the protocols that evaluate UPT. Gain = baseline $-$ UPT (positive favours UPT), paired 95\% interval clustered by network-window.}
\label{tab:all-endpoints}
\begin{tabular}{llrcc}
\toprule
Protocol & Endpoint & Gain & 95\% CI & Result \\
\midrule
\multicolumn{5}{l}{\emph{v6: METR-LA, PEMS-BAY, Intel}} \\
 & UPT vs transformer (PFN), NLL & +0.170 & [+0.136, +0.205] & pass \\
 & UPT non-inferior to 2-model pool & +0.126 & [+0.091, +0.161] & pass \\
 & UPT vs TabPFN v2, NLL, all sensors & -0.255 & [-0.315, -0.195] & \textbf{fail} \\
 & UPT vs PFN, synthetic, 16 sensors & +0.156 & [+0.136, +0.175] & pass \\
\multicolumn{5}{l}{\emph{v7: PeMS03, PeMS08, China air}} \\
 & UPT vs TabPFN v2, NLL, +3 dropped & +0.320 & [+0.302, +0.337] & pass \\
 & UPT vs TabPFN v2, NLL, +6 dropped & +0.989 & [+0.945, +1.033] & pass \\
 & UPT vs transformer (PFN), NLL & +0.361 & [+0.283, +0.439] & pass \\
 & float32 sites non-inferior to float64 & +0.000 & [-0.000, +0.000] & pass \\
 & (sec.) UPT-L vs TabPFN v2, all sensors & +0.313 & [+0.253, +0.372] & pass \\
 & (sec.) UPT-L vs TabPFN v2, +3 & +0.505 & [+0.450, +0.561] & pass \\
 & (sec.) UPT-L vs TabPFN v2, +6 & +1.073 & [+1.017, +1.128] & pass \\
 & (sec.) UPT-L vs UPT & +0.230 & [+0.174, +0.286] & pass \\
\multicolumn{5}{l}{\emph{v8: the six networks above}} \\
 & UPT vs TNP-D, NLL, all sensors & +0.382 & [+0.349, +0.415] & pass \\
 & UPT vs CNP, NLL, all sensors & +1.521 & [-0.252, +3.294] & \textbf{fail} \\
 & UPT vs calibrated kriging, NLL, all sensors & +1.014 & [+0.872, +1.156] & pass \\
 & UPT vs TNP-D, NLL, +3 dropped & +0.526 & [+0.487, +0.566] & pass \\
 & UPT vs CNP, NLL, +3 dropped & +0.782 & [+0.534, +1.029] & pass \\
 & UPT vs calibrated kriging, NLL, +3 dropped & +1.037 & [+0.889, +1.184] & pass \\
 & UPT vs TNP-D, NLL, +6 dropped & +1.046 & [+0.960, +1.132] & pass \\
 & UPT vs CNP, NLL, +6 dropped & +1.004 & [+0.740, +1.268] & pass \\
 & UPT vs calibrated kriging, NLL, +6 dropped & +1.324 & [+1.137, +1.511] & pass \\
\multicolumn{5}{l}{\emph{v9: DWD temp., AURN NO$_2$, USGS PA}} \\
 & UPT$^\ast$ vs TabPFN v2, NLL, all sensors & -0.374 & [-0.543, -0.205] & \textbf{fail} \\
 & UPT$^\ast$ vs TabPFN v2, NLL, +3 & -0.284 & [-0.440, -0.127] & \textbf{fail} \\
 & UPT$^\ast$ vs TabPFN v2, NLL, +6 & +0.077 & [-0.088, +0.243] & \textbf{fail} \\
 & UPT$^\ast$ vs PFN-L (matched scale), NLL & +0.029 & [+0.011, +0.047] & pass \\
 & UPT$^\ast$ vs TNP-D, NLL & +0.258 & [+0.194, +0.321] & pass \\
 & UPT$^\ast$ vs calibrated kriging, NLL & +0.708 & [+0.552, +0.863] & pass \\
 & UPT$^\ast$ vs TabPFN v2, CRPS$/\sigma$ & +0.012 & [+0.008, +0.016] & pass \\
\multicolumn{5}{l}{\emph{v10: USGS NY, USGS WA, DWD pressure, AURN PM$_{2.5}$}} \\
 & UPT-SF vs TabPFN v2, NLL, all sensors & -0.753 & [-0.927, -0.579] & \textbf{fail} \\
 & UPT-SF vs TabPFN v2, NLL, +3 & -0.631 & [-0.813, -0.449] & \textbf{fail} \\
 & UPT-SF vs TabPFN v2, NLL, +6 & -0.206 & [-0.442, +0.029] & \textbf{fail} \\
 & UPT-SF vs PFN-L (matched scale), NLL & -541.299 & [-1603.170, +520.572] & \textbf{fail} \\
 & UPT-SF vs TabPFN v2, CRPS$/\sigma$ & +0.002 & [-0.001, +0.005] & \textbf{fail} \\
 & UPT-SF vs TNP-D, NLL & -541.163 & [-1603.032, +520.706] & \textbf{fail} \\
 & UPT-SF vs calibrated kriging, NLL & -517.787 & [-1534.617, +499.043] & \textbf{fail} \\
\bottomrule
\end{tabular}
\end{table}

\section{Networks and episodes}
\label{app:data}

\begin{table}[h]\centering\scriptsize
\caption{\textbf{The thirteen sensor networks.} Test period year; protocol that introduced the network; test episodes and network-time-window clusters.}
\label{tab:networks}
\begin{tabular}{llllccc}
\toprule
Network & Quantity & Source & Test & Protocol & Episodes & Windows \\
\midrule
METR-LA & traffic speed & Los Angeles loop detectors~\citep{li2018dcrnn} & 2012 & v6 & -- & -- \\
PEMS-BAY & traffic speed & Bay Area loop detectors~\citep{li2018dcrnn} & 2017 & v6 & -- & -- \\
Intel Lab & indoor temperature & Intel Berkeley lab motes~\citep{intellab2004} & 2004 & v6 & -- & -- \\
PeMS03 & traffic flow & Sacramento detectors~\citep{song2020stsgcn} & 2018 & v7 & -- & -- \\
PeMS08 & traffic flow & San Bernardino detectors~\citep{guo2019astgcn} & 2016 & v7 & -- & -- \\
China air & PM$_{2.5}$ & 400 stations outside Beijing~\citep{zheng2015forecasting} & 2015 & v7 & -- & -- \\
DWD temp. & air temperature & German weather service stations~\citep{dwd2026} & 2026 & v9 & -- & -- \\
AURN NO$_2$ & NO$_2$ & UK air-quality network~\citep{carslaw2012openair} & 2026 & v9 & -- & -- \\
USGS PA & river discharge & Pennsylvania stream gauges~\citep{usgs2026nwis} & 2026 & v9 & -- & -- \\
USGS NY & river discharge & New York stream gauges~\citep{usgs2026nwis} & 2026 & v10 & -- & -- \\
USGS WA & river discharge & Washington stream gauges~\citep{usgs2026nwis} & 2026 & v10 & -- & -- \\
DWD pressure & air pressure & German weather service stations~\citep{dwd2026} & 2026 & v10 & -- & -- \\
AURN PM$_{2.5}$ & PM$_{2.5}$ & UK air-quality network~\citep{carslaw2012openair} & 2026 & v10 & -- & -- \\
\bottomrule
\end{tabular}
\end{table}

\paragraph{Construction.} Every network is processed identically:
\begin{enumerate}
  \item \textbf{Inputs:} the target sensor's $P=11$ inputs are its nearest sensors. For traffic flow and river discharge, they
   are the sensors most correlated with it over the source period.
  \item \textbf{Episodes:} each episode is one window of one target. It holds $48$ support and $48$ query rows, drawn from the
   rows where the target is observed.
  \item \textbf{Support dropout:} networks that are nearly complete receive $20\%$ simulated per-sensor dropout on the support
   rows.
  \item \textbf{Test targets:} these are a seeded sample of sensors.
  \item \textbf{Splits:} the source and test periods are consecutive and do not overlap.
\end{enumerate}

\paragraph{Snapshots.} The 2025--2026 data are provisional: publishers may revise them. We therefore record the SHA-256 of
each snapshot.

\section{Implementation}
\label{app:impl}

\paragraph{Architecture.} Small models use $d=64$, 4 layers, 4 heads and feed-forward width 128 (about 0.2M parameters). The
large models UPT-L and PFN-L use $d=128$, 8 layers, 8 heads and feed-forward width 512 (2.1M parameters).

\paragraph{Prior.} Synthetic tasks draw $P=5$ sensors from a one-factor model with nonlinear responses
$x_j=a_jy+b_j+0.4\sin(f_jy)+d_ju+s_j\epsilon_j$. Support missingness is random.

\paragraph{Pretraining.} AdamW with cosine decay.
\begin{itemize}
  \item Small models: 30k steps, 16 tasks $\times$ 16 queries, lr $10^{-3}$.
  \item Large models: 40k steps, 32 tasks $\times$ 16 queries, lr $5\times10^{-4}$.
\end{itemize}

\paragraph{Fine-tuning.} Identical for every model and network:
\begin{itemize}
  \item 4{,}000 source episodes and 2{,}000 AdamW steps;
  \item lr $3\times10^{-4}$ with cosine decay;
  \item 64 episodes $\times$ 8 queries per step;
  \item query masks are the natural mask with at most one extra dropped sensor.
\end{itemize}

\paragraph{Baselines.}
\begin{itemize}
  \item \textbf{TabPFN v2:} \texttt{tabpfn} package, ungated v2 weights, 8 estimators. Missing inputs are passed as NaN, and
  NLL is computed from its bar distribution in raw units.
  \item \textbf{TNP-D and CNP:} pretrained and fine-tuned with the same recipe as our models.
  \item \textbf{IGNNK-style kriging:} trained on the source period of each network. A per-episode affine recalibration on the
  support gives the calibrated variant.
  \item \textbf{Others:} LightGBM, MICE with Bayesian linear regression, and the closed forms (FA-Gaussian, EM-Gaussian, ridge
  and Bayesian linear regression), as defined in the repository.
\end{itemize}

\paragraph{Compute.}
\begin{itemize}
  \item Hardware: one A100 80GB, or an L4, on Colab.
  \item Fine-tuning UPT-L: about 15 minutes per network and seed.
  \item TabPFN v2: about 1 to 2 hours per network of 1{,}200--1{,}500 episodes $\times$ 3 conditions under shared load.
\end{itemize}

\section{Additional results}
\label{app:results}
\begin{table}[h]\centering\scriptsize
\caption{Per-network metrics, all sensors: NLL / CRPS$/\sigma$ / RMSE$/\sigma$ / 90\% coverage.}
\setlength{\tabcolsep}{2pt}
\begin{tabular}{lccccc}
\toprule
Network & UPT-SF & PFN-L & TabPFN v2 & TNP-D & Krig.\ (cal.) \\
\midrule
METR-LA & -1.82/0.183/0.561/0.91 & -1.48/0.194/0.569/0.94 & -1.79/--/0.615/-- & -1.37/0.212/0.586/0.94 & -0.50/0.229/0.569/0.90 \\
PEMS-BAY & -2.68/0.092/0.303/0.89 & -2.42/0.095/0.307/0.94 & -2.68/--/0.403/-- & -2.11/0.130/0.384/0.93 & -1.46/0.156/0.390/0.91 \\
Intel Lab & -0.95/0.058/0.259/0.87 & -0.32/0.060/0.198/0.95 & -1.30/--/0.177/-- & +0.03/0.088/0.257/0.93 & +0.94/0.241/0.533/0.90 \\
PeMS03 & -1.30/0.064/0.163/0.90 & -1.03/0.066/0.159/0.94 & -1.04/--/0.192/-- & -0.93/0.072/0.174/0.93 & -0.42/0.103/0.218/0.89 \\
PeMS08 & -2.54/0.087/0.236/0.92 & -1.12/0.092/0.231/0.94 & -1.82/--/0.281/-- & -1.07/0.097/0.238/0.94 & -0.68/0.125/0.295/0.86 \\
China air & -0.25/0.216/0.473/0.88 & -0.12/0.219/0.474/0.91 & -0.24/--/0.483/-- & -0.05/0.233/0.490/0.91 & +0.17/0.277/0.552/0.89 \\
DWD temp. & +1.02/0.101/0.205/0.89 & +1.04/0.101/0.204/0.90 & +1.32/0.130/0.254/0.87 & +1.19/0.114/0.226/0.89 & +1.39/0.135/0.261/0.90 \\
AURN NO$_2$ & +0.22/0.327/0.614/0.89 & +0.23/0.326/0.612/0.90 & +0.30/0.352/0.655/0.90 & +0.31/0.349/0.647/0.90 & +0.37/0.368/0.674/0.90 \\
USGS PA & -2.62/0.090/0.281/0.89 & -2.26/0.087/0.256/0.93 & -3.66/0.082/0.273/0.91 & -1.82/0.123/0.327/0.92 & -0.82/0.298/0.610/0.93 \\
USGS NY & +2024.70/0.144/0.490/0.88 & -1.56/0.143/0.466/0.92 & -3.63/0.131/0.450/0.91 & -1.41/0.180/0.516/0.91 & +85.31/0.320/0.693/0.93 \\
USGS WA & -3.31/0.099/0.266/0.89 & -2.73/0.097/0.250/0.91 & -4.97/0.091/0.250/0.92 & -2.51/0.124/0.309/0.91 & -1.95/0.238/0.516/0.89 \\
DWD pressure & -0.57/0.024/0.057/0.89 & -0.01/0.030/0.057/0.98 & -0.15/0.037/0.087/0.90 & +0.09/0.036/0.072/0.96 & +0.15/0.047/0.102/0.91 \\
AURN PM$_{2.5}$ & -0.40/0.246/0.530/0.89 & -0.27/0.248/0.528/0.91 & -0.43/0.270/0.564/0.88 & -0.21/0.260/0.543/0.91 & -0.07/0.285/0.572/0.90 \\
\bottomrule
\end{tabular}
\end{table}
\begin{table}[h]\centering\scriptsize
\caption{Per-network metrics, +3 sensors dropped: NLL / CRPS$/\sigma$ / RMSE$/\sigma$ / 90\% coverage.}
\setlength{\tabcolsep}{2pt}
\begin{tabular}{lccccc}
\toprule
Network & UPT-SF & PFN-L & TabPFN v2 & TNP-D & Krig.\ (cal.) \\
\midrule
METR-LA & -1.75/0.198/0.593/0.91 & -1.38/0.211/0.601/0.94 & -1.55/--/0.709/-- & -1.14/0.242/0.652/0.93 & -0.24/0.251/0.613/0.88 \\
PEMS-BAY & -2.56/0.105/0.338/0.90 & -2.30/0.108/0.339/0.94 & -2.46/--/0.466/-- & -1.83/0.157/0.450/0.92 & -1.25/0.171/0.424/0.89 \\
Intel Lab & -0.64/0.086/0.300/0.88 & +0.01/0.089/0.284/0.92 & +0.82/--/0.540/-- & +0.71/0.164/0.408/0.87 & +1.34/0.280/0.598/0.83 \\
PeMS03 & -1.18/0.071/0.175/0.91 & -0.94/0.072/0.172/0.94 & -0.81/--/0.238/-- & -0.75/0.086/0.199/0.92 & -0.43/0.104/0.222/0.90 \\
PeMS08 & -2.42/0.091/0.242/0.93 & -1.10/0.095/0.237/0.95 & -1.50/--/0.287/-- & -0.96/0.107/0.254/0.93 & -1.04/0.122/0.294/0.90 \\
China air & -0.17/0.232/0.499/0.89 & -0.05/0.234/0.499/0.91 & +0.23/--/0.629/-- & +0.04/0.257/0.529/0.91 & +0.28/0.299/0.588/0.87 \\
DWD temp. & +1.09/0.108/0.217/0.89 & +1.10/0.107/0.216/0.90 & +1.43/0.147/0.294/0.89 & +1.31/0.127/0.248/0.92 & +1.42/0.138/0.266/0.89 \\
AURN NO$_2$ & +0.29/0.350/0.654/0.88 & +0.31/0.349/0.652/0.89 & +0.40/0.386/0.712/0.87 & +0.40/0.379/0.697/0.89 & +0.43/0.387/0.706/0.88 \\
USGS PA & -2.50/0.101/0.309/0.90 & -2.16/0.097/0.279/0.93 & -3.36/0.116/0.352/0.92 & -1.54/0.156/0.393/0.90 & -0.70/0.330/0.660/0.89 \\
USGS NY & +1632.54/0.154/0.506/0.89 & -1.47/0.154/0.487/0.92 & -3.40/0.161/0.501/0.92 & -1.19/0.207/0.555/0.91 & +85.44/0.353/0.738/0.89 \\
USGS WA & -3.21/0.109/0.288/0.90 & -2.65/0.107/0.275/0.91 & -4.79/0.112/0.299/0.91 & -2.37/0.144/0.349/0.90 & -1.95/0.238/0.519/0.90 \\
DWD pressure & -0.44/0.027/0.064/0.89 & +0.03/0.033/0.064/0.97 & +0.21/0.058/0.149/0.94 & +0.31/0.048/0.100/0.92 & +0.27/0.053/0.117/0.88 \\
AURN PM$_{2.5}$ & -0.33/0.260/0.549/0.89 & -0.22/0.261/0.548/0.91 & -0.32/0.298/0.609/0.87 & -0.13/0.280/0.572/0.91 & -0.03/0.293/0.585/0.89 \\
\bottomrule
\end{tabular}
\end{table}
\begin{table}[h]\centering\scriptsize
\caption{Per-network metrics, +6 sensors dropped: NLL / CRPS$/\sigma$ / RMSE$/\sigma$ / 90\% coverage.}
\setlength{\tabcolsep}{2pt}
\begin{tabular}{lccccc}
\toprule
Network & UPT-SF & PFN-L & TabPFN v2 & TNP-D & Krig.\ (cal.) \\
\midrule
METR-LA & -1.60/0.230/0.669/0.90 & -1.09/0.248/0.674/0.94 & -1.21/--/0.827/-- & -0.44/0.298/0.760/0.89 & +0.30/0.292/0.711/0.84 \\
PEMS-BAY & -2.34/0.135/0.422/0.90 & -2.03/0.140/0.420/0.94 & -2.06/--/0.608/-- & -0.92/0.216/0.575/0.87 & -0.57/0.205/0.509/0.86 \\
Intel Lab & -0.02/0.150/0.436/0.85 & +0.79/0.157/0.425/0.86 & +2.36/--/0.813/-- & +1.75/0.292/0.609/0.78 & +2.72/0.357/0.716/0.73 \\
PeMS03 & -1.00/0.082/0.196/0.92 & -0.80/0.083/0.193/0.94 & -0.04/--/0.434/-- & -0.16/0.142/0.302/0.84 & -0.34/0.109/0.232/0.88 \\
PeMS08 & -2.12/0.097/0.255/0.92 & -1.01/0.101/0.248/0.95 & -0.62/--/0.460/-- & -0.43/0.142/0.306/0.87 & -0.94/0.130/0.304/0.87 \\
China air & -0.03/0.263/0.551/0.89 & +0.06/0.265/0.550/0.91 & +0.92/--/0.806/-- & +0.21/0.302/0.603/0.90 & +0.51/0.340/0.659/0.82 \\
DWD temp. & +1.22/0.121/0.240/0.89 & +1.25/0.119/0.237/0.88 & +2.14/0.265/0.495/0.72 & +1.62/0.168/0.317/0.91 & +1.49/0.145/0.278/0.87 \\
AURN NO$_2$ & +0.41/0.388/0.716/0.87 & +0.44/0.388/0.717/0.87 & +0.62/0.452/0.813/0.80 & +0.55/0.431/0.781/0.85 & +0.54/0.422/0.761/0.85 \\
USGS PA & -2.27/0.125/0.360/0.91 & -1.97/0.120/0.338/0.92 & -2.72/0.202/0.508/0.85 & -0.91/0.220/0.501/0.85 & -0.37/0.402/0.772/0.81 \\
USGS NY & +1305.58/0.176/0.537/0.89 & -1.34/0.179/0.527/0.91 & -2.88/0.236/0.609/0.86 & -0.75/0.258/0.622/0.87 & +85.80/0.436/0.859/0.81 \\
USGS WA & -2.98/0.128/0.333/0.90 & -2.47/0.129/0.328/0.90 & -4.41/0.154/0.385/0.86 & -2.11/0.177/0.412/0.88 & -1.78/0.274/0.588/0.85 \\
DWD pressure & -0.22/0.034/0.079/0.90 & +0.14/0.039/0.080/0.96 & +1.46/0.181/0.391/0.76 & +1.40/0.100/0.206/0.76 & +0.65/0.069/0.151/0.79 \\
AURN PM$_{2.5}$ & -0.21/0.285/0.587/0.89 & -0.10/0.286/0.586/0.90 & +0.07/0.383/0.728/0.77 & +0.06/0.328/0.645/0.88 & +0.08/0.315/0.617/0.86 \\
\bottomrule
\end{tabular}
\end{table}
